%% titre: Quiz 2 — Convexité et fonctionnelles quadratiques
%% matiere: Optimisation
%% niveau: Licence 3
%% annee: 2026-2027
%% ordre: 2
%% description: Ensembles et fonctions convexes, caractérisations différentielles, critère de Sylvester.

\begin{question}
Une fonction différentiable $f$ est convexe sur un convexe $C$ si et seulement si, pour tous $x,y\in C$ :
\begin{reponses}
  \bonne $f(y)\ge f(x)+\nabla f(x)^{\top}(y-x)$
  \mauvaise $f(y)\le f(x)+\nabla f(x)^{\top}(y-x)$
  \mauvaise $\nabla f(x)=0$
  \mauvaise $f(y)\ge f(x)$
\end{reponses}
\begin{explication}
Le graphe de $f$ est au-dessus de chacune de ses tangentes.
\end{explication}
\end{question}

\begin{question}
Lesquelles de ces fonctions sont convexes sur $\R$ ? (plusieurs réponses possibles)
\begin{reponses}
  \bonne $x\mapsto x^2$
  \bonne $x\mapsto x^4$
  \bonne $x\mapsto e^{x}$
  \mauvaise $x\mapsto x^3$
  \mauvaise $x\mapsto\sin x$
\end{reponses}
\begin{explication}
Pour une fonction de classe $\mathcal C^2$ sur $\R$, la convexité équivaut à $f''\ge0$ : $2$, $12x^2$ et $e^x$ sont positifs, alors que $(x^3)''=6x$ et $(\sin x)''=-\sin x$ changent de signe.
\end{explication}
\end{question}

\begin{question}
Laquelle de ces affirmations est vraie ?
\begin{reponses}
  \mauvaise Si $f$ est strictement convexe de classe $\mathcal C^2$, alors $\nabla^2f(x)$ est définie positive pour tout $x$.
  \bonne $x\mapsto x^4$ est strictement convexe bien que sa dérivée seconde s'annule en $0$.
  \mauvaise Si $f$ est convexe, alors $\nabla^2f(x)$ est définie positive pour tout $x$.
  \mauvaise Une hessienne semi-définie positive en un seul point suffit pour que $f$ soit convexe.
\end{reponses}
\begin{explication}
Hessienne définie positive partout $\Rightarrow$ strictement convexe, mais la réciproque est fausse : $(x^4)''(0)=0$.
\end{explication}
\end{question}

\begin{question}
Soit $f(x)=\frac12x^{\top}Ax-b^{\top}x+c$ avec $A$ symétrique. Son gradient est :
\begin{reponses}
  \mauvaise $Ax+b$
  \bonne $Ax-b$
  \mauvaise $\tfrac12Ax-b$
  \mauvaise $A$
\end{reponses}
\begin{explication}
On dérive composante par composante : $\partial_kf=(Ax)_k-b_k$ grâce à la symétrie de $A$. La hessienne est $A$ (constante).
\end{explication}
\end{question}

\begin{question}
Dans la situation précédente, $f$ est strictement convexe si et seulement si :
\begin{reponses}
  \mauvaise $A$ est inversible
  \mauvaise $\det A>0$
  \bonne $A$ est définie positive
  \mauvaise $b=0$
\end{reponses}
\begin{explication}
La hessienne $\nabla^2f=A$ est constante : $f$ est convexe si $A\succeq0$ et strictement convexe si $A\succ0$.
\end{explication}
\end{question}

\begin{question}
D'après le critère de Sylvester, une matrice symétrique $A$ est définie positive si et seulement si :
\begin{reponses}
  \mauvaise $\Delta_1>0$ seulement
  \mauvaise $\det A>0$
  \mauvaise tous les $\Delta_k\ge0$
  \bonne tous les mineurs principaux dominants $\Delta_k$ sont $>0$
\end{reponses}
\begin{explication}
$\Delta_k$ est le déterminant de la sous-matrice $k\times k$ du coin supérieur gauche. Avec des $\Delta_k$ nuls, on ne peut pas conclure pour le caractère semi-défini.
\end{explication}
\end{question}

\begin{question}
La matrice $A=\begin{pmatrix}1&2\\2&1\end{pmatrix}$ est :
\begin{reponses}
  \mauvaise définie positive
  \mauvaise définie négative
  \mauvaise semi-définie positive
  \bonne ni définie positive, ni définie négative
\end{reponses}
\begin{explication}
$\Delta_1=1>0$ mais $\Delta_2=1-4=-3<0$. Les valeurs propres sont $3$ et $-1$ : la matrice est indéfinie.
\end{explication}
\end{question}
